% !TEX program = xelatex
\documentclass[10pt]{beamer}
\usetheme{Madrid}
\usepackage{tikz} 
\usepackage{amsmath,amssymb,amsthm}

% 1. Hide navigation symbols
\setbeamertemplate{navigation symbols}{}

% 2. Define colors
\definecolor{MetroDark}{RGB}{26, 60, 115}   
\definecolor{MetroOrange}{RGB}{235, 129, 27} 
\definecolor{LightBg}{RGB}{240, 245, 250}

% 3. Apply colors
\setbeamercolor{palette primary}{bg=MetroDark,fg=white}
\setbeamercolor{palette secondary}{bg=MetroDark!80,fg=white}
\setbeamercolor{palette tertiary}{bg=MetroDark!90,fg=white}
\setbeamercolor{title}{bg=MetroDark,fg=white}
\setbeamercolor{frametitle}{bg=MetroDark,fg=white}
\setbeamercolor{itemize item}{fg=MetroOrange}
\setbeamercolor{itemize subitem}{fg=MetroOrange}
\setbeamercolor{structure}{fg=MetroDark}

\setbeamercolor{block title}{bg=MetroDark, fg=white}
\setbeamercolor{block body}{bg=LightBg, fg=black}
\setbeamercolor{block title alerted}{bg=MetroOrange, fg=white}
\setbeamercolor{block body alerted}{bg=MetroOrange!10, fg=black}

\setbeamertemplate{theorems}[numbered]
\setbeamertemplate{bibliography item}[text]
\newtheorem{proposition}[theorem]{Proposition} 

% --- Title Page ---
\title[Independent Edge Sampling for Spanners]{Random Edge Sampling \\ for Spanner Construction}
\subtitle{B.Sc. Thesis in Computer Science}
\author[M. Ahmadi]{Candidate: Mohammadreza Ahmadi \\ \vspace{0.3cm} Supervisor: Dr. Hossein Jowhari}
\institute[KNTU]{Faculty of Mathematics \\ K. N. Toosi University of Technology}
\date{September 2026}
\setbeamertemplate{footline}[frame number]

\begin{document}

\begin{frame}
    \titlepage
\end{frame}


% ==========================================
\section{Introduction and Preliminaries}
% ==========================================

\begin{frame}{Core Definitions: Graph Spanners}
    We consider unweighted, undirected $n$-vertex graphs $G = (V,E)$.
    
    \vspace{0.4cm}
    \begin{itemize}
        \item \textbf{Multiplicative $\alpha$-Spanner:} A subgraph $H = (V, E')$ preserving shortest paths up to a \textbf{stretch factor} $\alpha \ge 1$:
        \[ d_H(u,v) \le \alpha \cdot d_G(u,v) \quad \forall u, v \in V \]
        
        \vspace{0.6cm}
        \begin{proposition}
            $H$ is a $t$-spanner of $G \iff d_H(u, v) \le t$ for all edges $uv \in E$.
        \end{proposition}
    \end{itemize}
\end{frame}

\begin{frame}{Example: A $2$-Spanner}
    \begin{center}
        \begin{tikzpicture}[scale=1.2]
            \tikzstyle{every node}=[circle, draw=MetroDark, fill=MetroDark, inner sep=0pt, minimum size=4pt]
            
            \begin{scope}[xshift=-3.5cm]
                \node (v1) at (0, 0) {};
                \node (v2) at (1.5, 0.5) {};
                \node (v3) at (2.5, 0) {};
                \node (v4) at (0.5, 1.5) {};
                \node (v5) at (2, 2) {};
                \node (v6) at (3, 1.5) {};
                
                \draw[gray, thick] (v1) -- (v2) -- (v3) -- (v6) -- (v5) -- (v4) -- (v1);
                \draw[gray, thick] (v2) -- (v5);
                \draw[gray, thick] (v2) -- (v4);
                \draw[gray, thick] (v3) -- (v5);
                \draw[gray, thick] (v1) -- (v5);
                \node[draw=none, fill=none] at (1.5, -0.6) {Original Graph $G$};
            \end{scope}

            \draw[->, very thick, MetroOrange] (-0.3, 1) -- (0.5, 1);

            \begin{scope}[xshift=1.5cm]
                \node (u1) at (0, 0) {};
                \node (u2) at (1.5, 0.5) {};
                \node (u3) at (2.5, 0) {};
                \node (u4) at (0.5, 1.5) {};
                \node (u5) at (2, 2) {};
                \node (u6) at (3, 1.5) {};
                
                \draw[MetroDark, thick] (u1) -- (u2) -- (u3);
                \draw[MetroDark, thick] (u3) -- (u6) -- (u5) -- (u4) -- (u1);
                \draw[MetroDark, thick] (u2) -- (u5);
                \node[draw=none, fill=none] at (1.5, -0.6) {A $2$-Spanner $H$};
            \end{scope}
        \end{tikzpicture}
    \end{center}
    \vspace{0.3cm}
    \begin{itemize}
        \item \textbf{Intractability:} Finding the optimally sparse $t$-spanner is \textbf{NP-Hard}, forcing us to use approximations.
    \end{itemize}
\end{frame}

\begin{frame}{Theoretical Limits: Girth and Edge Density}
    \textbf{Girth ($g$):} Length of the shortest cycle in a graph.
    
    \vspace{0.2cm}
    \begin{proposition}
        A graph with girth $> t + 1$ has \textbf{no proper subgraph} that is a $t$-spanner.
    \end{proposition}
    
    \vspace{0.4cm}
    \textbf{Maximum Edge Density ($\gamma$):}
    \begin{itemize}
        \item $\gamma(n,k)$: Max possible edges in an $n$-vertex graph with girth $> k$.
        \item \textbf{The Fundamental Barrier:} Because of the proposition above, no algorithm can construct a $t$-spanner sparser than $\gamma(n, t+1)$ in the worst case.
    \end{itemize}
\end{frame}

\begin{frame}{The Asymptotic Baseline: Moore \& Erdős}
    Exact values for $\gamma(n,k)$ remain a major open problem, but we rely on asymptotic bounds:
    
    \vspace{0.4cm}
    \begin{block}{Moore Bound \& Erdős' Girth Conjecture}
        For graphs with girth $> 2k$:
        \begin{itemize}
            \item \textbf{Upper Bound (Moore):} $O(n^{1+1/k})$ edges.
            \item \textbf{Lower Bound (Erdős):} $\Omega(n^{1+1/k})$ edges (Conjectured tight).
        \end{itemize}
    \end{block}
    
    \vspace{0.4cm}
    \begin{itemize}
        \item A $(2k-1)$-spanner with $O(n^{1+1/k})$ edges is considered structurally \textbf{optimal}.
    \end{itemize}
\end{frame}

\begin{frame}{The Classical Approach: Greedy Algorithm}
    \begin{block}{Greedy Spanner (Althöfer et al., 1993)}
        Process edges sequentially. Add $e=(u,v)$ to $H$ \textbf{only if}:
        \[ d_H(u,v) > 2k-1 \]
    \end{block}
    
    \vspace{0.6cm}
    \begin{center}
        \renewcommand{\arraystretch}{1.5} % فاصله‌گذاری مناسب بین ردیف‌ها
        \begin{tabular}{p{0.45\textwidth} | p{0.45\textwidth}}
            \textcolor{MetroDark}{\Large \textbf{Pros}} & \textcolor{MetroOrange}{\Large \textbf{Cons}} \\
            \hline
            \textcolor{MetroDark}{$\checkmark$} \textbf{Optimal Sparsity} \newline {\footnotesize Guarantees $O(n^{1+1/k})$ edges.} 
            & \textcolor{MetroOrange}{$\times$} \textbf{High Time Complexity} \newline {\footnotesize $O(k \cdot n^{2+1/k})$ bottleneck.} \\
            
            & \textcolor{MetroOrange}{$\times$} \textbf{Strictly Sequential} \newline {\footnotesize Fails in distributed networks.}
        \end{tabular}
    \end{center}
\end{frame}

% ==========================================
\section{State of the Art \& Our Motivation}
% ==========================================

\begin{frame}{State of the Art: Elkin \& Neiman (2019)}
    \textbf{Goal:} Bypass the Greedy algorithm's sequential bottleneck via a distributed randomized framework.
    
    \vspace{0.4cm}
    \textbf{The Mechanism: \textcolor{MetroOrange}{Vertex-Based Randomness}}
    \begin{itemize}
        \item Each vertex $u \in V$ independently samples $r_u \sim \mathcal{EXP}(\ln(cn)/k)$.
        \item Vertices broadcast these radii to neighbors and select edges based on the maximums received.
    \end{itemize}

    \vspace{0.4cm}
    \textbf{Theoretical Achievements:}
    \begin{itemize}
        \item Achieves $O((n/\epsilon)^{1+1/k})$ expected edges.
        \item Computes the spanner in exactly $k$ communication rounds (CONGEST model).
    \end{itemize}
\end{frame}

\begin{frame}{The Paradigm Shift: Our Motivation}
    The Elkin-Neiman framework is elegant, but inherently algorithmic:
    
    \vspace{0.3cm}
    \begin{itemize}
        \item Randomness is exclusively assigned to the \textbf{vertices}.
        \item It requires message passing, broadcasting, and state comparisons.
    \end{itemize}

    \vspace{0.8cm}
    \begin{alertblock}{The Core Question of Our Thesis}
        Can we eliminate these communication phases entirely? 
        
        \vspace{0.3cm}
        \textit{What if we sample edges \textbf{independently and directly}, relying solely on the local topology of the graph?}
    \end{alertblock}
\end{frame}

\begin{frame}{Independent Edge Sampling: Initial Attempts}
    How do we define the independent sampling probability $p_e$?
    
    \vspace{0.4cm}
    \textbf{Attempt 1: Global Minimum Probability}
    \begin{itemize}
        \item Set a uniform, high probability across all edges to ensure path survival.
        \item \textcolor{MetroOrange}{Result:} The graph remains far too dense (fails the sparsity goal).
    \end{itemize}

    \vspace{0.6cm}
    \textbf{Attempt 2: Deterministic Deletion}
    \begin{itemize}
        \item Drop an edge $e$ simply because it has an alternative path $P$.
        \item \textcolor{MetroOrange}{Result:} Edges inside $P$ drop independently (cascading failures). The path collapses, destroying the stretch guarantee.
    \end{itemize}
\end{frame}

\begin{frame}{The Dependency Bottleneck}
    \textbf{Definition:} Let $x_e$ be the maximum number of edge-disjoint \textbf{good paths} (length $\le t$) alternative to $e$.
    
    \vspace{0.1cm}
    \[
    p_e = 
    \begin{cases} 
        1 & \text{if } x_e \le c \\
        \frac{c}{x_e} & \text{if } x_e > c
    \end{cases}
    \]
    
    \vspace{0.3cm}
    \textbf{Result:}
    \begin{itemize}
        \item Path survival requires \textit{every} internal edge to be sampled.
        \item A single dense edge ($x_{e'} \gg c$) in a good path has $p_{e'} \to 0$.
        \item This single dropped edge disconnects the path and violates the stretch.
    \end{itemize}
\end{frame}

% ==========================================
\section{The $2c$-Cutoff Framework}
% ==========================================

\begin{frame}{Our Contribution: The $2c$-Cutoff Filter}
    \textbf{Core Idea:} Alternative paths must be structurally \textbf{reliable}.
    
    \vspace{0.3cm}
    \begin{block}{Definition: Safe Path}
        A path $P$ is a \textbf{safe path} for edge $e$ if:
        \begin{enumerate}
            \item Edge-disjoint from all other safe paths for $e$.
            \item Length $\le t = 2k-1$.
            \item Contains no highly dense edges: $x_{e'} \le 2c \;\; (\forall e' \in P)$.
        \end{enumerate}
    \end{block}

    \vspace{0.4cm}
    To ensure path survival, the cutoff is set to:
    \[
\boldsymbol{c = \Theta\!\left(2^t \log n\right)}
\]
\end{frame}

\begin{frame}{Expected Sparsity Analysis}
    \textbf{Goal:} Bound the expected number of edges $\mathbb{E}[|E'|] = \sum p_e$.
    
    \vspace{0.4cm}
    \textbf{1. Critical Edges ($x_e \le c$) $\implies p_e = 1$}
    \begin{itemize}
        \item Bounded via Greedy Decomposition and Moore's Bound.
        \item \textbf{Total size:} $\mathbf{O(c \cdot n^{1+1/k})}$.
    \end{itemize}

    \vspace{0.6cm}
    \textbf{2. Dense Edges ($x_e > c$) $\implies p_e = c/x_e$}
    \begin{itemize}
        \item Partitioned into $O(\log n)$ exponentially growing buckets.
        \item Bucket growth rate perfectly cancels the exponential drop in $p_e$.
        \item \textbf{Total size:} $\mathbf{O(c \cdot n^{1+1/k} \log n)}$.
    \end{itemize}
\end{frame}

\begin{frame}{Final Spanner Size \& Performance}
    Substituting $c = \Theta(2^{2k} \log n)$, the expected number of edges is:
    \[ \mathbf{O(2^{2k} n^{1+1/k} \log^2 n)} \]
    
    \vspace{0.6cm}
    \textbf{Constant Stretch Regime:}
    \begin{itemize}
        \item For small constant $k$, the $2^{2k}$ overhead is a trivial factor.
        \item The framework naturally produces highly sparse spanners purely through independent sampling.
    \end{itemize}
    
    \vspace{0.4cm}
    \textbf{Optimal Scaling:}
    \begin{itemize}
        \item By optimizing the scale parameter to $k^* = \Theta(\sqrt{\log n})$.
        \item The total expected size of the spanner becomes near-linear: $\mathbf{n^{1+o(1)}}$.
    \end{itemize}
\end{frame}

% ==========================================
\section{Algorithmic Realization}
% ==========================================

\begin{frame}{The Algorithmic Challenge}
    \begin{alertblock}{The NP-Hardness}
        Computing $x_e$ (the exact maximum number of edge-disjoint safe paths) is \textbf{NP-Hard} for $t \ge 4$ \textcolor{gray}{\small (Itai et al., 1982)}.
    \end{alertblock}
    \vspace{1cm}
    We need an efficiently computable parameter $q_e$ that can safely replace $x_e$ without violating the stretch guarantees.
\end{frame}

\begin{frame}{The Two-Phase Greedy Algorithm}    
    \begin{block}{Phase 1: Upper Bound Extraction}
        \begin{itemize}
            \item Temporarily \textbf{ignore} the safety condition.
            \item Run iterative BFS to extract a maximal set of paths. Let this count be $q_e^{(0)}$.
            \item If $q_e^{(0)} \le 2c$, the edge is definitively labeled \textbf{safe}. Otherwise, it is \textbf{heavy}.
        \end{itemize}
    \end{block}

    \vspace{0.3cm}
    \begin{block}{Phase 2: Safe Path Isolation}
        \begin{itemize}
            \item Remove all \textit{heavy} edges from the graph.
            \item Re-run iterative BFS on this purified subgraph.
            \item The new count is $q_e$. Every path found is structurally guaranteed to be valid and safe.
        \end{itemize}
    \end{block}
    
    \vspace{0.2cm}
    \textbf{Time Complexity:} Iterative BFS takes $O(m)$ per path $\implies \mathbf{O(m^3)}$ overall.
\end{frame}

\begin{frame}{Algorithmic Realization: Greedy Approximation}
    \begin{block}{Lemma: Greedy $t$-Approximation}
        The extracted maximal set size ($q_e$) bounds the optimal maximum set size ($x_e$):
        \[ \frac{x_e}{t} \le q_e \le x_e \]
    \end{block}
    
    \vspace{0.4cm}
    \textbf{Proof:}
    \begin{itemize}
        \item Let $OPT$ ($|OPT| = x_e$) be the \textit{maximum} possible packing of safe paths.
        \item Let $GREEDY$ ($|GREEDY| = q_e$) be our extracted \textit{maximal} packing.
    \end{itemize}
    
    \vspace{0.4cm}
    \textbf{The Upper Bound ($q_e \le x_e$)}
    \begin{itemize}
        \item Trivial by definition. Since $x_e$ is the absolute theoretical maximum capacity, no valid packing ($GREEDY$) can physically exceed it.
    \end{itemize}
\end{frame}

\begin{frame}{Proof Continued: The Lower Bound ($x_e \le t \cdot q_e$)}
    \textbf{The Lower Bound}
    
    \vspace{0.3cm}
    \begin{itemize}
        \item Since $GREEDY$ is \textit{maximal}, it cannot be extended. Thus, every path in $OPT$ must intersect at least one path in $GREEDY$.
        
        \vspace{0.2cm}
        \item A single path $P \in GREEDY$ has at most $t$ edges.
        
        \vspace{0.2cm}
        \item Because paths in $OPT$ are strictly edge-disjoint, each edge of $P$ can intersect at most \textit{one} path in $OPT$.
        
        \vspace{0.2cm}
        \item Consequently, a single path in $GREEDY$ can block no more than $t$ distinct paths in $OPT$.
    \end{itemize}
    
    \vspace{0.4cm}
    \begin{block}{Conclusion}
        Summing the maximum blocked paths across all of $GREEDY$:
        \[ |OPT| \le t \cdot |GREEDY| \implies \mathbf{x_e \le t \cdot q_e} \hfill \square \]
    \end{block}
\end{frame}

\begin{frame}{The Final Algorithmic Spanner Size}
    Since $q_e \ge x_e/t$, the approximation increases the sampling probability by at most a factor of $t$.
    
    \vspace{0.5cm}
    \textbf{Linearity of Expectation:}
    \[ 
    \mathbb{E}[|E'|] = \sum_{e \in E} p'_e \le t \sum_{e \in E} p_e = t \cdot \mathbb{E}[|E_{opt}|] 
    \]
    
    \vspace{0.6cm}
    \begin{exampleblock}{Polynomial-Time Guarantee}
        Substituting $t = 2k-1$, the final expected spanner size becomes:
        \[ \mathbf{O(k \cdot 2^{2k} n^{1+1/k} \log^2 n)} \]
    \end{exampleblock}
\end{frame}

% ==========================================
\section{Conclusion & Future Work}
% ==========================================

\begin{frame}{Conclusion \& Open Problems}
    \textbf{Key Contributions:}
    \begin{itemize}
        \item Designed a novel framework based purely on independent edge sampling.
        \item Resolved the dependency bottleneck using the $2c$-cutoff filter.
        \item Bypassed NP-Hardness via an $O(m^3)$ two-phase greedy approximation.
        \item Achieved an expected spanner size of $O(k \cdot 2^{2k} n^{1+1/k} \log^2 n)$.
    \end{itemize}

    \vspace{0.6cm}
    \textbf{Future Directions:}
    \begin{itemize}
        \item Is the $O(2^{2k})$ overhead fundamental to independent sampling, or merely an artifact of our analysis?
        \item Can the $O(m^3)$ time complexity or the $t$-approximation ratio be further improved?
        \item Are there entirely different local topological metrics that avoid path-counting ($x_e$) altogether?
    \end{itemize}
\end{frame}

\begin{frame}
    \vspace{1cm}
    \begin{center}
        \Huge \textbf{Thank You!} \\
    \end{center}
\end{frame}

\end{document}